\documentclass[11pt,a4paper]{article}
\usepackage[utf8]{inputenc}
\usepackage[T1]{fontenc}
\renewcommand{\contentsname}{Table des matières}
\usepackage{amsmath,amssymb}
\usepackage{geometry}
\geometry{margin=2.2cm}
\usepackage{xcolor}
\usepackage{enumitem}
\usepackage{booktabs}
\usepackage{tabularx}
\usepackage{mdframed}
\usepackage{hyperref}
\hypersetup{colorlinks=true,linkcolor=blue!50!black}

\definecolor{cbleu}{RGB}{25,80,180}
\definecolor{cvert}{RGB}{20,120,55}
\definecolor{cgris}{RGB}{95,95,95}

\newmdenv[
  linecolor=cgris!60, backgroundcolor=cgris!5,
  linewidth=0.8pt, roundcorner=3pt, innermargin=8pt, innertopmargin=6pt, innerbottommargin=6pt,
  skipabove=8pt, skipbelow=10pt
]{ex}

\newmdenv[
  linecolor=cbleu, backgroundcolor=cbleu!5,
  linewidth=1pt, roundcorner=4pt, innermargin=8pt, innertopmargin=6pt, innerbottommargin=6pt,
  skipabove=8pt, skipbelow=8pt
]{info}

\newcounter{exo}
\newcommand{\exo}[1]{\refstepcounter{exo}\par\medskip\noindent
{\color{cbleu}\bfseries\large Exercice \theexo\ --- #1}\par\smallskip}

\title{\textbf{Exercices d'entraînement}\\[3pt]
\large Énoncés seuls -- coups de pouce -- réponses finales}
\author{}
\date{}

\begin{document}
\maketitle
\vspace{-0.8cm}

\begin{info}
\textbf{Comment utiliser ce document.} Les corrections détaillées, vous les avez déjà ailleurs : ici, le
but est de vous \textbf{tester à blanc}. Faites chaque exercice entièrement seul, avec uniquement le
formulaire. Si vous bloquez, allez d'abord voir le \textbf{coup de pouce} (partie 2) -- il donne la
direction sans donner le calcul. N'allez aux \textbf{réponses} (partie 3) qu'une fois l'exercice terminé.

\smallskip
Données communes : $R=8{,}314$ J/(mol$\cdot$K) ; air $\gamma=1{,}4$, $21\ \%$ d'O$_2$, $M=28{,}9$ g/mol ;
$c_{eau}=4185$ J/(kg$\cdot$K).
\end{info}

\section{Énoncés}

\exo{Moteur essence 4 temps}
\begin{ex}
Un moteur à essence 4 temps possède \textbf{6 cylindres} et une cylindrée totale de \textbf{1,8 L}. Son
rapport volumétrique vaut $\varepsilon=9{,}5$. L'air est aspiré à $20\,^\circ$C sous 1 bar
($\gamma=1{,}4$). Le carburant a pour formule simplifiée \textbf{CH$_{1,9}$}, l'excès d'air vaut
$\lambda=1{,}1$ et son PCI vaut \textbf{43 000 kJ/kg}. Les rendements valent $\eta_{forme}=0{,}72$ et
$\eta_{m\acute{e}ca}=0{,}82$. Le moteur tourne à \textbf{3000 tr/min}.

\begin{enumerate}[label=\alph*),leftmargin=1.6em,nosep]
\item Calculez $V_A$, $V_B$ et le nombre de moles admises par cylindre.
\item Déterminez $T_C$ et $p_C$ en fin de compression.
\item Calculez la masse de carburant brûlée par cycle et par cylindre, puis $Q_{CD}$.
\item Quel est le rendement théorique du cycle ?
\item Déduisez-en le travail disponible et la \textbf{puissance effective} du moteur.
\end{enumerate}
\end{ex}

\exo{Moteur diesel 4 temps}
\begin{ex}
Un moteur diesel 4 temps a \textbf{6 cylindres} et une cylindrée totale de \textbf{3 L}. Son rapport
volumétrique vaut $\varepsilon=17$ et son rapport de combustion $\rho=2{,}2$. L'air est admis à
$25\,^\circ$C sous 1 bar ($\gamma=1{,}4$). Rendements : $\eta_{forme}=0{,}78$,
$\eta_{m\acute{e}ca}=0{,}85$. Régime : \textbf{1800 tr/min}.

\begin{enumerate}[label=\alph*),leftmargin=1.6em,nosep]
\item Donnez $V_A$, $V_B$, $V_D$ et le nombre de moles par cylindre.
\item Calculez $T$ et $p$ aux points C, D et E.
\item Calculez la chaleur $Q_{CD}$ apportée par la combustion.
\item Calculez le rendement théorique, puis la puissance effective.
\item Pourquoi le rendement théorique d'un diesel est-il supérieur à celui de l'essence de
l'exercice 1, alors que la formule est plus « défavorable » ?
\end{enumerate}
\end{ex}

\exo{Moteur 2 temps}
\begin{ex}
Un petit moteur 2 temps \textbf{monocylindre} a un rapport volumétrique $\varepsilon=8{,}5$
($\gamma=1{,}4$). Chaque cycle apporte $Q_{CD}=95$ J. Les rendements valent $\eta_{forme}=0{,}68$ et
$\eta_{m\acute{e}ca}=0{,}80$. Il tourne à \textbf{5000 tr/min}.

\begin{enumerate}[label=\alph*),leftmargin=1.6em,nosep]
\item Calculez le rendement théorique.
\item Calculez le travail disponible par cycle.
\item Calculez la puissance effective.
\item Quelle puissance obtiendrait-on si ce même moteur était un 4 temps tournant au même régime ?
Commentez.
\end{enumerate}
\end{ex}

\exo{Machine frigorifique}
\begin{ex}
Une machine frigorifique au R134a doit produire une puissance frigorifique de \textbf{5 kW}. Lecture sur
le diagramme des frigoristes : $h_1=395$ kJ/kg (sortie évaporateur), $h_{2s}=428$ kJ/kg (point
isentropique), $h_3=250$ kJ/kg (sortie condenseur). Le rendement isentropique du compresseur vaut
$\eta_{isos}=0{,}7$.

\begin{enumerate}[label=\alph*),leftmargin=1.6em,nosep]
\item Calculez l'enthalpie réelle $h_2$ en sortie de compresseur, et donnez $h_4$.
\item Calculez $w$, $q_1$ et $q_2$ par kilogramme de fluide, et vérifiez votre bilan.
\item Calculez la \textbf{PSN} de cette machine.
\item Déterminez le débit de fluide frigorigène et la puissance du compresseur.
\item Quelle puissance le condenseur évacue-t-il ? Vérifiez la cohérence.
\end{enumerate}
\end{ex}

\exo{Pompe à chaleur avec circuit d'eau}
\begin{ex}
Une pompe à chaleur assure le chauffage d'une maison qui perd \textbf{9 kW}. L'évaporateur travaille à
$-8\,^\circ$C (air extérieur) et le condenseur à $40\,^\circ$C. Sur le diagramme : $h_1=400$,
$h_2=440$, $h_3=256$ kJ/kg. Le condenseur réchauffe l'eau d'un circuit de radiateurs de
\textbf{$6\,^\circ$C}.

\begin{enumerate}[label=\alph*),leftmargin=1.6em,nosep]
\item Calculez $w$, $q_1$, $q_2$ et vérifiez le bilan du cycle.
\item Calculez le COP de cette pompe à chaleur.
\item Calculez le COP maximal théorique entre ces deux sources. Le COP réel est-il crédible ?
\item Déterminez le débit de fluide frigorigène et la puissance électrique du compresseur.
\item Calculez le débit d'eau du circuit de chauffage, en kg/s puis en m$^3$/h.
\end{enumerate}
\end{ex}

\exo{Turbine à gaz (cycle de Joule)}
\begin{ex}
Une turbine à gaz fonctionne à l'air ($c_p=1004$ J/kg$\cdot$K, $\gamma=1{,}4$) avec un débit massique de
\textbf{20 kg/s}. L'air entre dans le compresseur à $20\,^\circ$C sous 1 bar et en ressort à
\textbf{8 bar} (compression adiabatique réversible). Il est ensuite chauffé de façon isobare jusqu'à
\textbf{$900\,^\circ$C} dans la chambre de combustion, puis détendu de façon adiabatique réversible
jusqu'à 1 bar dans la turbine. Un refroidissement isobare ferme le cycle.

\begin{enumerate}[label=\alph*),leftmargin=1.6em,nosep]
\item Calculez $T_2$ (sortie compresseur) et $T_4$ (sortie turbine).
\item Calculez les travaux $w_{12}$ et $w_{34}$ par kg, puis le travail net.
\item Calculez la chaleur $q_{23}$ fournie dans la chambre de combustion.
\item Déduisez-en le rendement du cycle. Retrouvez-le par la formule $\eta=1-r_p^{-(\gamma-1)/\gamma}$
où $r_p$ est le rapport de pression : les deux résultats coïncident-ils ?
\item Calculez la puissance nette de l'installation.
\end{enumerate}
\end{ex}

\newpage
\section{Coups de pouce}
\emph{À ne lire que si vous bloquez -- ils indiquent la direction, pas le calcul.}

\begin{center}
\renewcommand{\arraystretch}{1.5}
\begin{tabularx}{\linewidth}{@{}c X@{}}
\toprule
\textbf{Ex.} & \textbf{Indication} \\
\midrule
\textbf{1} & Le système $\{V_B-V_A=\text{cyl. unitaire}\ ;\ V_B/V_A=\varepsilon\}$ donne
$V_A=\frac{\text{cyl}}{\varepsilon-1}$. Pour c), n'oubliez pas que les $n$ moles du cylindre contiennent
\emph{air + carburant} : $n=n_{carb}(n_{air/carb}+1)$. \\
\textbf{2} & La combustion est \textbf{isobare} : $T_D/T_C=V_D/V_C=\rho$ directement, et $Q_{CD}$ se
calcule avec $C_p$ (pas $C_v$). Pour la détente, le rapport de volumes est $\rho/\varepsilon$. \\
\textbf{3} & Aucune donnée de combustion n'est nécessaire : $\eta_{th}$ ne dépend que de $\varepsilon$ et
$\gamma$. Attention à la formule de puissance : en 2 temps, 1 tour $=$ 1 cycle. \\
\textbf{4} & $\eta_{isos}=\frac{h_{2s}-h_1}{h_2-h_1}$ (le réel est \emph{plus grand} que l'isentropique).
Pour un \textbf{frigo}, l'utile est $q_2$ (évaporateur) : c'est lui qui fixe le débit. \\
\textbf{5} & Pour une \textbf{PAC}, l'utile est $|q_1|$ (condenseur) : c'est lui qui fixe le débit. Le
circuit d'eau se traite avec $P=\dot m\,c\,\Delta T$, la puissance étant la même des deux côtés de
l'échangeur. \\
\textbf{6} & Machines \textbf{ouvertes} : travaillez par kg avec $c_p$, pas avec des moles. Pour les
adiabatiques, utilisez la forme reliant $T$ et $p$ : $T_2=T_1\,(p_2/p_1)^{(\gamma-1)/\gamma}$. Ici les
deux isobares sont à $p=8$ bar et $p=1$ bar. \\
\bottomrule
\end{tabularx}
\end{center}

\newpage
\section{Réponses finales}
\emph{Seulement les valeurs, pour que vous puissiez vous corriger sans voir la méthode. Un écart de
quelques dixièmes est normal (arrondis).}

\subsection*{Exercice 1 -- moteur essence}
\begin{center}
\begin{tabularx}{0.9\linewidth}{@{}X r@{}}
\toprule
$V_A=V_C$ & $35{,}29$ cm$^3$ \\
$V_B$ & $335{,}29$ cm$^3$ \\
$n$ (par cylindre) & $0{,}01376$ mol \\
$T_C$ \ / \ $p_C$ & $721{,}4$ K \ / \ $23{,}38$ bar \\
air réel par mole de carburant & $7{,}726$ mol \\
$m_{carb}$ par cycle et cylindre & $2{,}191\cdot10^{-5}$ kg $=21{,}9$ mg \\
$Q_{CD}$ & $942{,}3$ J \\
$\eta_{th}$ & $\mathbf{59{,}36\ \%}$ \\
$|W_{th}|$ \ / \ $|W_{ind}|$ \ / \ $|W_{dispo}|$ & $559{,}4$ \ / \ $402{,}8$ \ / \ $330{,}3$ J \\
$\eta_{eff}$ & $35{,}05\ \%$ \\
$P_{eff}$ & $\mathbf{49{,}5\ kW}$ \\
\midrule
\emph{(pour vérifier)} $T_D$ \ / \ $p_D$ \ / \ $T_E$ & $4017$ K \ / \ $130{,}2$ bar \ / \ $1632$ K \\
\bottomrule
\end{tabularx}
\end{center}

\subsection*{Exercice 2 -- moteur diesel}
\begin{center}
\begin{tabularx}{0.9\linewidth}{@{}X r@{}}
\toprule
$V_A=V_C$ \ / \ $V_B$ \ / \ $V_D$ & $31{,}25$ \ / \ $531{,}25$ \ / \ $68{,}75$ cm$^3$ \\
$n$ (par cylindre) & $0{,}02143$ mol \\
$T_C$ \ / \ $p_C$ & $926{,}0$ K \ / \ $52{,}80$ bar \\
$T_D$ \ / \ $p_D$ & $2037{,}2$ K \ / \ $52{,}80$ bar \\
$T_E$ \ / \ $p_E$ & $899{,}1$ K \ / \ $3{,}02$ bar \\
$Q_{CD}$ & $693{,}0$ J \\
$\eta_{th}$ & $\mathbf{61{,}37\ \%}$ \\
$|W_{th}|$ \ / \ $|W_{dispo}|$ & $425{,}3$ \ / \ $282{,}0$ J \\
$P_{eff}$ & $\mathbf{25{,}4\ kW}$ \\
\bottomrule
\end{tabularx}
\end{center}
\textbf{e)} Parce que le rendement dépend surtout de $\varepsilon$, et que le diesel comprime beaucoup
plus fort ($17$ contre $9{,}5$). Le terme en $\rho$ pénalise effectivement le diesel à $\varepsilon$ égal,
mais ce handicap est largement compensé par le rapport volumétrique bien plus élevé -- que le diesel peut
se permettre justement parce qu'il ne comprime que de l'air (pas de risque d'auto-allumage prématuré).

\subsection*{Exercice 3 -- moteur 2 temps}
\begin{center}
\begin{tabularx}{0.9\linewidth}{@{}X r@{}}
\toprule
$\eta_{th}$ & $\mathbf{57{,}52\ \%}$ \\
$|W_{th}|$ \ / \ $|W_{ind}|$ \ / \ $|W_{dispo}|$ & $54{,}64$ \ / \ $37{,}15$ \ / \ $29{,}72$ J \\
$P_{eff}$ (2 temps) & $\mathbf{2{,}48\ kW}$ \\
$P_{eff}$ si c'était un 4 temps & $1{,}24$ kW \\
\bottomrule
\end{tabularx}
\end{center}
\textbf{d)} Exactement \textbf{deux fois moins} : à régime égal, un 4 temps ne réalise qu'un cycle tous
les deux tours. C'est l'avantage principal du 2 temps (plus de puissance à cylindrée égale), payé par un
rendement réel plus faible (une partie du mélange frais part directement à l'échappement).

\subsection*{Exercice 4 -- machine frigorifique}
\begin{center}
\begin{tabularx}{0.9\linewidth}{@{}X r@{}}
\toprule
$h_2$ réel \ / \ $h_4$ & $442{,}14$ \ / \ $250$ kJ/kg \\
$w$ \ / \ $q_1$ \ / \ $q_2$ & $+47{,}14$ \ / \ $-192{,}14$ \ / \ $+145{,}0$ kJ/kg \\
vérification $q_1+q_2+w$ & $0$ \checkmark \\
\textbf{PSN} & $\mathbf{3{,}08}$ \\
débit de fluide & $\mathbf{0{,}0345\ kg/s}$ \\
$P_{compresseur}$ & $\mathbf{1{,}63\ kW}$ \\
$P_{condenseur}$ & $6{,}63$ kW $=5+1{,}63$ \checkmark \\
\bottomrule
\end{tabularx}
\end{center}

\subsection*{Exercice 5 -- pompe à chaleur}
\begin{center}
\begin{tabularx}{0.9\linewidth}{@{}X r@{}}
\toprule
$w$ \ / \ $q_1$ \ / \ $q_2$ & $+40$ \ / \ $-184$ \ / \ $+144$ kJ/kg \\
vérification & $0$ \checkmark \\
\textbf{COP} & $\mathbf{4{,}60}$ \\
$COP_{max}$ (Carnot) & $6{,}52$ \ -- le réel vaut $71\ \%$ du max : crédible \checkmark \\
débit de fluide & $\mathbf{0{,}0489\ kg/s}$ \\
$P_{compresseur}$ & $\mathbf{1{,}96\ kW}$ \\
débit d'eau & $\mathbf{0{,}358\ kg/s}=1{,}29$ m$^3$/h \\
\bottomrule
\end{tabularx}
\end{center}

\subsection*{Exercice 6 -- turbine à gaz}
\begin{center}
\begin{tabularx}{0.9\linewidth}{@{}X r@{}}
\toprule
$T_2$ & $531{,}0$ K \\
$T_4$ & $647{,}6$ K \\
$w_{12}$ (compresseur) & $+238{,}8$ kJ/kg \\
$w_{34}$ (turbine) & $-527{,}6$ kJ/kg \\
$w_{net}$ & $-288{,}8$ kJ/kg \\
$q_{23}$ & $+644{,}7$ kJ/kg \\
$\eta$ par les travaux & $\mathbf{44{,}8\ \%}$ \\
$\eta$ par la formule $1-r_p^{-(\gamma-1)/\gamma}=1-8^{-0,2857}$ & $\mathbf{44{,}8\ \%}$ \checkmark \\
$P_{nette}$ & $\mathbf{-5{,}78\ MW}$ \\
\bottomrule
\end{tabularx}
\end{center}
\textbf{d)} Oui, les deux méthodes donnent le même résultat : c'est la meilleure vérification possible.
Notez au passage que pour un cycle de Joule \emph{simple} (une seule turbine, pas de reprise), le
rendement ne dépend que du \textbf{rapport de pression}, exactement comme le rendement d'un moteur à
essence ne dépend que de $\varepsilon$.

\end{document}
